% =========================================================
% ABSTRACT
% El encabezado (\chapter*) está en main.tex.
% =========================================================

Insecurity in Mexico exposes citizens to latent risks during their daily commutes. However, pedestrians often lack analytical tools to evaluate the level of vulnerability associated with their chosen routes. This research proposes the development of a mobile application prototype designed to generate safe pedestrian routes, grounded in optimal pathfinding algorithms (Shortest Path Problem), Graph Theory, crowdmapping techniques, and Geographic Information Systems (GIS). The computational model weighs variables related to crime incidence, urban morphology, Environmental Criminology, and citizen perception in order to compute paths with minimal risk exposure.

\vspace{0.5cm}
\noindent\textbf{Keywords:} Mobile application, Graph Theory, Geographic Information Systems, Environmental Criminology.
